\subsection{The complete dyadic locus on the cyclic-eight genus-two cover}

\begin{statement}{A.192}{Proved}
{Four compact Selmer classes and the single cyclic-eight character}
For \(C,J,D,T\) in (A.189.1), let
\(j_C(R)=[R-\infty]\).  The complete local intersection,
with its global coordinate, is
\[
\begin{aligned}
 \{(\eta,R)\in\operatorname{Sel}(\Q,T_2J)\times C(\Q_2):
       \operatorname{loc}_2\eta=\kappa_2(j_C(R))\}
  =\{&(0,\infty),(\kappa_2(T),(0,0)),\\
     &(\kappa_2(D),(3,9)),(\kappa_2(-D),(3,-9))\}.
\end{aligned}
\tag{A.192.1}
\]
In particular
\(C(\Q)=\{\infty,(0,0),(3,9),(3,-9)\}\).

There is an actual connected degree-eight cover of the tripod
\[
 W:\quad z^8=-\frac{t(1-t)^3}{24},\qquad t(1-t)z\ne0,
\tag{A.192.2}
\]
whose proper compactification is \(C\).  Every locally
point-or-cusp-matching full metabelian tripod section with
\[
                   \kappa_8\bigl(t(1-t)^3\bigr)=[-24]
\tag{A.192.3}
\]
has at two either an ordinary centre exactly three, a cusp
at zero, or a cusp at one.  Its selected compact cover
image is respectively \(D\) or \(-D\), \(T\), or zero
in \(\operatorname{Sel}(\Q,T_2J)\).  The same assertion
holds for full Selmer sections in this finite class.
Only these local and compact-coefficient assertions are made.
\end{statement}

\begin{proof}
Retain the actual logarithm and its slope \(\rho\equiv1\pmod4\)
from A.189.  The linear functional
\(\Lambda(L_1,L_2)=L_2-\rho L_1\) vanishes on the
entire global completed Mordell--Weil image.  A.189 treats
all points with \(x\) a unit: the square condition on
\(x^5-54x\) modulo eight forces \(x\equiv3\pmod8\).
We compute \(\Lambda\) on every remaining local point.

\emph{The exact comparison across the two positive-valuation reductions.}
The coordinates \(x=2X,y=2Y\) give the smooth integral
model
\[
                 \mathcal X:\quad Y^2=8X^5-27X.
\tag{A.192.4}
\]
Its special fibre is the affine line with coordinate \(Y\),
since \(X=Y^2\) there.  Its derivative with respect to
\(X\) is a unit on the special fibre.  The Abel--Jacobi
map \(a:R\mapsto[R-(0,0)]\) extends to the Neron
model \(\mathcal J\) by the actual smooth-source
mapping property used in A.189.

We prove that the special-fibre map \([4]a\) is zero.
Work first over an algebraic closure of the residue field.
The image of the connected affine line belongs to
\(\mathcal J_s^0\).  Its projection to the abelian
quotient of this connected algebraic group is constant.
Indeed a map from the affine line to an abelian variety
extends to the projective line by properness.  A nonconstant
map from the projective line would, after factoring a
finite power of Frobenius, have nonzero differential;
pulling back a suitable invariant differential would
give a nonzero regular differential on the projective
line.  This is impossible.  The Frobenius factorization
terminates because a nonconstant rational function has
finite inseparable degree.  Thus the image is in the
connected affine part.  Its projection to the torus
quotient is also constant, since every invertible
regular function on the affine line is constant.
Both constants are the identity because \(a(0,0)=0\).
The image therefore belongs to the smooth connected
unipotent radical \(U\), whose dimension is at most two.
These actual subgroup and quotient maps are the
perfect-field chain in
\cite[\S5.9.1, equation (5.9.1), p.147]{LegPoonenRationalPoints}.

Here \([4]U=0\), with a direct dimension proof.
A unitriangular embedding exists by
\cite[Theorem 5.4.8(ii), p.128]{LegPoonenRationalPoints}.
It makes a sufficiently high power of \([2]\) zero.
Take the reduced connected images of successive
\([2]\)-maps; over the perfect residue field these
are smooth connected unipotent subgroups.  A nonzero
image cannot have the same dimension as its predecessor:
it would then be the whole predecessor, making \([2]\)
surjective there, contrary to its nilpotence.  Each
nonzero step thus lowers dimension.  Starting with
dimension at most two proves \([4]U=0\).
The argument proves a morphism identity, since the
source is reduced; it descends to the residue field.
Consequently \([4]a\) has identically zero special-fibre
map on the entire affine line.

To make its analytic consequence explicit, let
\(A(Z)\in\Z[[Z]]\) be the unique solution of
\(A=1+ZA^5\), with constant term one.  Its coefficients
are integral by the recursive equation.  The formula
\[
        X=-\frac{Y^2}{27}
                    A\left(\frac{8Y^8}{27^5}\right)
\tag{A.192.5}
\]
is a restricted integral power series, convergent on
all \(Y\in\Z_2\).  It parametrizes the integral
points of (A.192.4): the equation
\(X=-Y^2/27+(8/27)X^5\) is a strict contraction
on \(\Z_2\) for each \(Y\).
After this substitution, the formal coordinates of
\([4]a\) belong to \(2\Z_2\langle Y\rangle\),
because its special-fibre map is zero.  The formal
logarithm therefore converges as a restricted analytic
function on this whole closed disc.  Explicitly,
after a fixed scalar multiplication the invariant
forms have integral coefficients, and their formal
primitive coefficients have denominators dividing
their total degrees.  Substituting elements of
\(2\Z_2\langle Y\rangle\) gives uniform convergence.
Its derivative is four times the pulled-back invariant
form, and its value at \(Y=0\) is zero.  Two restricted
power series in characteristic zero with the same
derivative and constant agree.  Division by four
therefore identifies the primitives below with the
actual logarithm of \([R-(0,0)]\), including when
\(Y\) is a unit.  This is the required comparison
across the two distinct residue points; convergence
of an inverse coordinate series alone is not used.

\emph{All positive valuations.}
Let \(r=v_2(x)>0\).  In \(x^5-54x\) the second
term has smaller valuation, so \(r=2k-1\) and
\(v_2(y)=k\geq1\).  Set \(z=y\).  Equation
(A.192.5) becomes
\[
       x=-\frac{z^2}{54}
                       A\left(\frac{z^8}{54^5}\right).
\tag{A.192.6}
\]
Write \(A(Z)=\sum a_nZ^n\) and
\(A(Z)^2=\sum b_nZ^n\); all coefficients are
integers and \(a_0=b_0=1\).  Differentiating and
integrating the actual forms \(dx/y,x\,dx/y\)
gives the logarithm relative to the torsion point:
\[
\begin{aligned}
 L_1&=-\sum_{n\geq0}
       \frac{(8n+2)a_nz^{8n+1}}
                 {(8n+1)54^{5n+1}},\\
 L_2&=\sum_{n\geq0}
       \frac{(4n+2)b_nz^{8n+3}}
                 {(8n+3)54^{5n+2}}.
\end{aligned}
\tag{A.192.7}
\]
Their leading terms are \(-z/27\) and
\(z^3/4374\).  Each later term has valuation
at least, respectively,
\(k+n(8k-5)\) and
\(3k-1+n(8k-5)\).
Since \(8k-5\geq3\), all tails converge and are
strictly higher than their leading terms.  Therefore
\[
                  v_2(L_1)=k,\qquad v_2(L_2)=3k-1,
                  \qquad\Lambda(L_1,L_2)\ne0.
\tag{A.192.8}
\]
The logarithm of \(T\) is zero, so these are also
the logarithm coordinates of \([R-\infty]\).

\emph{All negative valuations and the first infinity disc.}
If \(r=v_2(x)<0\), the term \(x^5\) dominates,
so \(r=-2k\), \(k\geq1\).  Put
\(w=1/x\), \(t_\infty=-x^2/y\); then
\[
        w=t_\infty^2(1-54w^4),\qquad
        x=1/w,\qquad y=-1/(w^2t_\infty).
\tag{A.192.9}
\]
This is the actual smooth integral infinity chart:
the derivative with respect to \(w\) is a unit
at \((w,t_\infty)=(0,0)\).  Here
\(v_2(t_\infty)=k\geq1\), so every point in
question belongs to that single formal residue
disc, including \(k=1\).
Write \(w=t_\infty^2B(t_\infty^8)\), where
\(B(Z)=1-54ZB(Z)^4\).
The recursive coefficients of \(B\) are integral,
and every nonconstant coefficient is divisible
by fifty-four.  Direct substitution in the
differentials gives
\[
       \omega_1=t_\infty\,dw,\qquad
       \omega_2=t_\infty\,dw/w.
\tag{A.192.10}
\]
Consequently the two primitives at infinity have
leading terms \(2t_\infty^3/3\) and
\(2t_\infty\).
For the first, the remaining coefficients of
\(t_\infty^{8n+3}\), \(n\geq1\), have
valuation at least two: differentiating contributes
\(8n+2\), and the integration denominator is odd.
For the second use
\(B'/B\in54\Z_2[[Z]]\); the coefficients of
\(t_\infty^{8n+1}\), \(n\geq1\), have
valuation at least four, because differentiation
contributes a factor eight.  These estimates prove
convergence and
\[
                v_2(L_1)=3k+1,\qquad v_2(L_2)=k+1,
                \qquad\Lambda(L_1,L_2)\ne0.
\tag{A.192.11}
\]
The same formal-group comparison as A.189 applies
on the explicitly smooth infinity chart.

The positive, negative, and unit cases exhaust
all affine points except \((0,0)\).  The latter,
\(\infty\), and \((3,\pm9)\) are actual rational
points.  A.186 and A.187 identify the global
two-adic Selmer group with \(\Z_2D\oplus(\Z/2)T\).
Its localization is injective by the nonzero
logarithm of \(D\) and injection of its finite
two-primary torsion, as proved in A.189.
The preceding exclusion and this injectivity
prove exactly (A.192.1), and hence the rational
point assertion as well.

\emph{The smaller cyclic cover and the section maps.}
The two function-field maps are literal inverses:
\[
 \begin{aligned}
 t&=\frac{x^4}{x^4-54},& z&=\frac{3x}{y},\\
 x&=-\frac{6z^2}{1-t},& y&=-\frac{18z}{1-t}.
 \end{aligned}
\tag{A.192.12}
\]
Substitution verifies both equations and both
compositions.  The eighth-root equation is
geometrically irreducible, since its radicand
has valuation one at \(t=0\); it gives a
degree-eight finite etale cover of \(U\).
Its proper compactification is therefore \(C\).
The boundary fibres are \((0,0)\) above zero,
\(\infty\) above one, and the four points
\((x,0)\) with \(x^4=54\) above infinity.
Their ramification indices are respectively
eight, eight, and two.  The last fibre has no
dyadic rational point, by valuation.
The two remaining rational points map to
\(t=81/(81-54)=3\).

Let \(H\) be the geometric subgroup of this
cover; it is the kernel of the actual character
\(\Pi^{\rm ab}\to\Z/8(1)\),
\((a,b)\mapsto a+3b\).
Equation (A.192.3) makes the pullback eighth-root
torsor trivial on the section.  Selecting its
fixed sheet and conjugating the finite section
accordingly gives the arithmetic cover extension.
As \(\Pi'\subset H\),
\(\Pi''\subset\overline{[H,H]}\).
Thus the full metabelian section has an actual
image in \(\pi_1(W)/\overline{[H,H]}\).
Compactification and \(j_C\) give its map to
the \(T_2J\) arithmetic extension.
The ordinary fixed-sheet and cuspidal
decomposition-group comparisons of A.178--A.189
give the actual local Kummer centres at every
place, including the real component.  Every
finite class is therefore in the corresponding
global Selmer group, and their compatible
limit is in \(\operatorname{Sel}(\Q,T_2J)\).
Applying (A.192.1) at two gives exactly the
four proper centres and the claimed coefficient
classes.  A proper centre on the boundary
comes from a cuspidal image; an ordinary
point of \(U\) lifts to its finite etale
fibre in \(W\), so it can only give one
of \((3,\pm9)\).  This proves the stated
tripod conclusions with their exact domains.
\end{proof}
